% Generated draft: requires mathematical and source review.
\documentclass[11pt]{article}
\usepackage[margin=25mm]{geometry}
\usepackage{fontspec}
\setmainfont{Latin Modern Roman}
\usepackage{amsmath,amssymb,mathtools,unicode-math}
\setmathfont{Latin Modern Math}
\usepackage{enumitem,needspace,xurl,hyperref}
\setlength{\parindent}{0pt}
\setlength{\parskip}{5pt}
\setlength{\emergencystretch}{3em}
\newcommand{\N}{\mathbb{N}}
\newcommand{\Z}{\mathbb{Z}}
\newcommand{\Q}{\mathbb{Q}}
\newcommand{\R}{\mathbb{R}}
\newcommand{\C}{\mathbb{C}}
\newcommand{\F}{\mathbb{F}}
\newcommand{\A}{\mathbb{A}}
\newcommand{\PP}{\mathbb{P}}
\newcommand{\E}{\mathbb{E}}
\newcommand{\eps}{\varepsilon}
\newcommand{\dd}{\,\mathrm{d}}
\begin{document}
\section*{UnsolvedMath ID: 30003076}
\textbf{Source problem code:} catalogue entry
\subsection{Definitions and mathematical statement}
Write \((X,H)\) for the smooth polarized surface in the supplied record, \(v\) for the prescribed sheaf class or Chern character, and \(M_H(v)\) for the indicated moduli space of sheaves with invariant \(v\). The source does not specify the base field, the precise meaning of polarization, the permitted sheaves, or the stability condition defining \(M_H(v)\). These choices must be fixed before a fully specified mathematical problem results; no particular choices are asserted here.

Once these choices are fixed in a setting with finite-dimensional sheaf cohomology, write
\[
 h^i(X,E)=\dim_k H^i(X,E),
 \qquad
 \chi(E)=\sum_{i\geq 0}(-1)^i h^i(X,E),
\]
where \(k\) is the chosen base field and the Euler characteristic is defined when the sum is finite. The notation \(\chi(v)\) presupposes that this Euler characteristic is determined by \(v\). In the usual setting of coherent sheaves on a smooth projective surface, only degrees \(0,1,2\) can contribute.

For a nonempty irreducible component \(C\subseteq M_H(v)\), interpret a property of a general sheaf in \(C\) as a property holding throughout some nonempty Zariski-open subset of \(C\). The supplied record does not specify a component or assert that \(M_H(v)\) is nonempty or irreducible, so the intended scope of the word ``general'' remains unresolved.

\textbf{Problem.} Determine the general cohomology of \(E\) in \(M_H(v)\): after resolving the preceding choices, determine the generic values of \(h^i(X,E)\) on the intended component or components. In particular, characterize when there is a nonempty Zariski-open subset \(U\subseteq C\) such that
\[
 \#\{i\geq 0:H^i(X,E)\neq 0\}\leq 1
 \qquad\text{for every }[E]\in U.
\]
If \(\chi(v)=0\), determine whether a general sheaf is acyclic, namely whether there is such a subset \(U\) with
\[
 H^i(X,E)=0
 \qquad\text{for every }[E]\in U\text{ and every }i\geq 0.
\]
These are questions about the applicable surfaces and sheaf classes, not assertions of universal vanishing. The source leaves unresolved whether the final question is intended independently or within the cases satisfying the preceding one-degree condition.
\subsection{Short English statement}
Determine the cohomology of a general sheaf with prescribed invariant on a smooth polarized surface. Characterize when its cohomology is concentrated in at most one degree, and ask whether Euler characteristic zero gives vanishing in every degree. The base field, sheaf and stability conventions, and intended moduli components require clarification.
\subsection{Sources}
\begin{itemize}
\item Supplied local registry record, permanent UnsolvedMath ID \(30003076\), source problem code \texttt{catalogue entry}, title \emph{Cohomology of General Sheaves on Polarized Surfaces}. The project identifies \texttt{data/unsolvedmath/problems.json} and \texttt{catalog/problems.json} as local registries; the supplied record does not identify which file contains this entry.
\item \emph{Mini-Workshop: Arrangements of Subvarieties, and their Applications in Algebraic Geometry} (2016), Oberwolfach Reports, pp.\ 653--702, DOI \texttt{10.4171/owr/2016/14}: \url{https://doi.org/10.4171/owr/2016/14}. Also listed in the supplied literature sources as ``OWR report 14/2016.'' The record attributes the problem to the workshop report and supplies no individual proposer.
\item Supplied literature reference labeled ``Coskun--Huizenga'': \url{https://arxiv.org/abs/1611.02674}. No title or precise theorem reference is supplied.
\item Supplied literature reference labeled ``Coskun--Huizenga--Raha'': \url{https://doi.org/10.1093/imrn/rnaf064}. No title or precise theorem reference is supplied.
\item Literature assessment dated \texttt{2026-08-21}: the supplied curation labels the problem ``partially solved'' and states that weak Brill--Noether determines general cohomology in substantial cases, including the projective plane and classified families on Hirzebruch surfaces, while the question for general polarized surfaces and Chern characters remains open. This is supplied dated curation derived from the OWR record, not an independently verified assessment. Its time-sensitive status and cited claims require independent verification.
\end{itemize}
\end{document}
